% Generated draft: requires mathematical and source review.
\documentclass[11pt]{article}
\usepackage[margin=25mm]{geometry}
\usepackage{fontspec}
\setmainfont{Latin Modern Roman}
\usepackage{amsmath,amssymb,mathtools,unicode-math}
\setmathfont{Latin Modern Math}
\usepackage{enumitem,needspace,xurl,hyperref}
\setlength{\parindent}{0pt}
\setlength{\parskip}{5pt}
\setlength{\emergencystretch}{3em}
\newcommand{\N}{\mathbb{N}}
\newcommand{\Z}{\mathbb{Z}}
\newcommand{\Q}{\mathbb{Q}}
\newcommand{\R}{\mathbb{R}}
\newcommand{\C}{\mathbb{C}}
\newcommand{\F}{\mathbb{F}}
\newcommand{\A}{\mathbb{A}}
\newcommand{\PP}{\mathbb{P}}
\newcommand{\E}{\mathbb{E}}
\newcommand{\eps}{\varepsilon}
\newcommand{\dd}{\,\mathrm{d}}
\begin{document}
\section*{UnsolvedMath ID: 9400179}
\textbf{Source problem code:} catalogue entry
\subsection{Definitions and mathematical statement}
Fix a prime $p$ and a totally real number field $F$. Inside a fixed algebraic closure let $\widetilde F$ be the compositum of all Galois $\Z_p$-extensions of $F$, and put $\Gamma=\operatorname{Gal}(\widetilde F/F)$. Let $\widetilde L$ be the maximal unramified abelian pro-$p$ extension of $\widetilde F$ and set
\[
\widetilde X=\operatorname{Gal}(\widetilde L/\widetilde F),\qquad
\widetilde\Lambda=\Z_p[[\Gamma]].
\]
The completed group algebra is the inverse limit of $\Z_p[\Gamma/U]$ over open subgroups $U$. Conjugation by lifts of $\Gamma$ gives the natural module action on the abelian group $\widetilde X$, independent of the chosen lifts.

In the standard finitely generated module convention over a commutative noetherian Iwasawa algebra, pseudo-null means that the support has codimension at least two; equivalently
\[
(\widetilde X)_{\mathfrak p}=0
\quad\text{for every prime }\mathfrak p\subset\widetilde\Lambda
\text{ of height at most one}.
\]
The question is whether $\widetilde X$ has this property. It is stronger than an unspecified assertion of torsion and is not the assertion that $\widetilde X=0$.

The source specifically restricts $F$ to be totally real; no extension to all number fields is inserted from the generalized-conjecture label. No rank-one assumption on $\Gamma$, or identification of the compositum with just the cyclotomic extension, is added. The usual structural/finiteness conventions needed for the module/support definition should be checked in the intended Iwasawa setting. The source also leaves unramifiedness at infinite places unspecified, which matters especially at $p=2$.
\subsection{Short English statement}
Is the unramified abelian pro-$p$ Galois module over the compositum of all $\Z_p$-extensions of a totally real field supported only in codimension at least two over its Iwasawa algebra?
\subsection{Sources}
\begin{itemize}
\item Local UnsolvedMath registry, permanent ID 9400179, exact source code catalogue entry, ``Greenberg's pseudo-null conjecture''; attributed to Ralph Greenberg, no proposed year supplied. Registry status open, classification OPEN-TRIAGE.
\item Wikipedia contributors, \emph{List of unsolved problems in mathematics}, Algebraic number theory bullet 50 and linked-page generalized-conjecture resolution, pinned revision 1366636767, \url{https://en.wikipedia.org/w/index.php?title=List_of_unsolved_problems_in_mathematics&oldid=1366636767#Number_theory}. Supplied MediaWiki-plus-linked-page extraction accessed July 29, 2026. The linked-page URL and a primary conjecture citation are not supplied here.
\item The research summary reports special abelian/cyclotomic and other family cases while calling generality open, but gives no precise theorem or hypotheses. No broad special-case claim is adopted as a theorem for the exact totally-real formulation.
\item Provenance flags current primary-literature status and redistribution terms NEEDS\_REVIEW; the family link was decomposed into distinct conjectures rather than treated as one record.
\end{itemize}
\end{document}
